\documentclass[../m00-session-02.tex]{subfiles}
\begin{document}
\TopicStandaloneTitle{03}{Finding Axes of Maximum Truth: Eigen \& SVD}{Applied Linear Algebra via NumPy}
\TopicDivider{03}{Finding Axes of Maximum Truth: Eigen \& SVD}{Invariant directions under transformation, spectral theorem, and PCA from scratch.}

% -----------------------------------------------------------------------------
% Frame 1: Cognitive Objectives
% -----------------------------------------------------------------------------
\begin{frame}{Cognitive Objectives}
  \begin{LearningObjectives}{By the end of this topic, learners will be able to:}
    \begin{itemize}\setlength{\itemsep}{0.28em}
      \item \textbf{Identify invariant directions} geometrically where linear transformations act as pure scalars ($\mathbf{A}\mathbf{v} = \lambda\mathbf{v}$).
      \item \textbf{Diagonalize symmetric operators} using the Spectral Theorem on data covariance matrices $\Sigma = \frac{1}{n}\mathbf{X}^T\mathbf{X}$.
      \item \textbf{Deconstruct SVD geometrically} as a 3-step pipeline: Rotation ($V^T$) $\to$ Scaling ($\Sigma$) $\to$ Rotation ($U$).
      \item \textbf{Implement PCA from scratch} using Singular Value Decomposition to compress high-dimensional feature spaces.
    \end{itemize}
  \end{LearningObjectives}
\end{frame}

% -----------------------------------------------------------------------------
% Frame 2: The Socratic Hook (The 50-Feature Redundancy Dilemma)
% -----------------------------------------------------------------------------
\begin{frame}[fragile]{The High-Dimensional Redundancy Dilemma}
  \begin{columns}[T,onlytextwidth]
    \begin{column}{0.48\textwidth}
      {\small\bfseries\color{UpGradRed}The 50-Feature Multicollinearity Curse}\par\vspace{0.08cm}
      \begin{itemize}\setlength{\itemsep}{0.22em}\small
        \item Customer transaction dataset contains $d = 50$ numeric columns.
        \item Many features (e.g., monthly spend, annual income, credit limit) are heavily correlated.
        \item \textcolor{UpGradRed}{\textbf{Curse of Dimensionality:}} Models overfit, distance metrics degrade, latency escalates.
        \item \textbf{Challenge:} How do we extract the top 3 orthogonal directions in space that preserve $95\%$ of all data variance?
      \end{itemize}
    \end{column}
    \hfill{\color{UpGradRule}\vrule width .5pt}\hfill
    \begin{column}{0.48\textwidth}
      {\small\bfseries\color{AWSEmerald}The Discovery: Invariant Axes of Variance}\par\vspace{0.08cm}
      \begin{itemize}\setlength{\itemsep}{0.22em}\small
        \item Under transformation, most vectors rotate out of their original line of action.
        \item \textbf{Eigenvectors} never rotate; they only stretch by factor $\lambda$:
          \[
          \mathbf{A}\mathbf{v} = \lambda\mathbf{v}
          \]
        \item On the data covariance matrix, the top eigenvector points along the line of \textbf{maximum spread / variance}!
      \end{itemize}
      \vspace{0.06cm}
      \TakeawayBanner{Eigenvectors find the natural coordinate system of the data.}
    \end{column}
  \end{columns}
\end{frame}

% -----------------------------------------------------------------------------
% Frame 3: Hero Visualization: Eigenvectors as Invariant Directions
% -----------------------------------------------------------------------------
\begin{frame}{Hero Visualization: Eigenvectors as Invariant Directions}
  \VisualExploration[.54]{%
    \begin{tikzpicture}[scale=0.92,>=stealth,every node/.style={transform shape}]
      % Axes
      \draw[->, thick, UpGradSlate!60] (-0.4, 0) -- (5.2, 0) node[right, font=\tiny\sffamily] {$x_1$};
      \draw[->, thick, UpGradSlate!60] (0, -0.4) -- (0, 3.8) node[above, font=\tiny\sffamily] {$x_2$};
      
      % Eigenvector span line (Dashed)
      \draw[dashed, line width=1.1pt, AWSBlue!70] (-0.3, -0.2) -- (4.8, 3.2) node[above, font=\scriptsize\bfseries\color{AWSBlue}] {$\text{Span}(\mathbf{v}_1)$};
      
      % Eigenvector before transformation
      \coordinate (V) at (1.8, 1.2);
      \draw[->, line width=2.2pt, AWSBlue] (0, 0) -- (V) node[below right=-1pt, font=\small\bfseries] {$\mathbf{v}_1$};
      
      % Eigenvector after transformation (Scaled by lambda=2.2 along same line!)
      \coordinate (AV) at (3.96, 2.64);
      \draw[->, line width=2.4pt, UpGradRed] (0, 0) -- (AV) node[above left, font=\small\bfseries] {$\mathbf{A}\mathbf{v}_1 = \lambda_1 \mathbf{v}_1$};
      
      % Non-eigenvector x (arbitrary direction)
      \coordinate (X) at (1.0, 2.5);
      \draw[->, line width=1.5pt, UpGradSlate] (0, 0) -- (X) node[left, font=\small\bfseries] {$\mathbf{x}$};
      
      % Non-eigenvector Ax (Rotates and scales!)
      \coordinate (AX) at (3.4, 0.8);
      \draw[->, line width=1.6pt, AWSOrange] (0, 0) -- (AX) node[below right, font=\small\bfseries] {$\mathbf{A}\mathbf{x}$};
      
      % Rotation arc for x
      \draw[->, thick, AWSOrange] (0.7, 1.75) arc (68.2:13.2:0.75);
      \node[font=\scriptsize\bfseries\color{AWSOrange}] at (1.4, 1.5) {Rotates off span!};
    \end{tikzpicture}%
  }{Invariant Directions Under $\mathbf{A}$}{%
    \begin{itemize}\setlength{\itemsep}{0.25em}
      \item \textbf{General Vector $\mathbf{x}$:} Multiplied by $\mathbf{A}$, changes both length \textbf{and} direction ($\mathbf{A}\mathbf{x} \nparallel \mathbf{x}$).
      \item \textbf{Eigenvector $\mathbf{v}$:} Multiplied by $\mathbf{A}$, remains on its original span line ($\mathbf{A}\mathbf{v} \parallel \mathbf{v}$).
      \item \textbf{Eigenvalue $\lambda$:} The scalar stretch factor along the invariant line.
      \item \textbf{Spectral Theorem:} Any symmetric matrix $\mathbf{S} = \mathbf{S}^T$ has an orthonormal basis of eigenvectors ($\mathbf{S} = \mathbf{Q}\Lambda\mathbf{Q}^T$).
    \end{itemize}
  }
\end{frame}

% -----------------------------------------------------------------------------
% Frame 4: Hero Visualization: The 3-Step SVD Transformation Pipeline
% -----------------------------------------------------------------------------
\begin{frame}{Hero Visualization: SVD Geometric Decomposition}
  \begin{center}
    \begin{tikzpicture}[scale=0.82,every node/.style={transform shape}]
      % Step 1: Unit Circle in R^n (v1, v2 basis)
      \begin{scope}[shift={(0,0)}]
        \node[font=\small\bfseries\color{UpGradNight}] at (0, 2.3) {1. Input Space ($\mathbf{V}^T$)};
        \draw[thick, UpGradRule!70] (-1.8,0) -- (1.8,0);
        \draw[thick, UpGradRule!70] (0,-1.8) -- (0,1.8);
        \draw[thick, AWSBlue, fill=AWSBlue!12] (0,0) circle (1.2cm);
        \draw[->, line width=1.8pt, UpGradRed] (0,0) -- (1.2,0) node[above right=-2pt, font=\small\bfseries] {$\mathbf{v}_1$};
        \draw[->, line width=1.8pt, AWSBlue] (0,0) -- (0,1.2) node[above left=-2pt, font=\small\bfseries] {$\mathbf{v}_2$};
      \end{scope}

      % Arrow 1: Rotate / Align
      \draw[->, ultra thick, UpGradSlate] (2.0, 0) -- (3.2, 0) node[midway, above=2pt, font=\small\bfseries\color{UpGradNight}] {$\mathbf{V}^T$};

      % Step 2: Singular Value Stretch (Sigma)
      \begin{scope}[shift={(5.2,0)}]
        \node[font=\small\bfseries\color{UpGradNight}] at (0, 2.3) {2. Axis Scaling ($\Sigma$)};
        \draw[thick, UpGradRule!70] (-2.2,0) -- (2.2,0);
        \draw[thick, UpGradRule!70] (0,-1.8) -- (0,1.8);
        \draw[thick, AWSOrange, fill=AWSOrange!12] (0,0) ellipse [x radius=1.9cm, y radius=0.7cm];
        \draw[->, line width=1.8pt, UpGradRed] (0,0) -- (1.9,0) node[below=2pt, font=\small\bfseries] {$\sigma_1 \mathbf{e}_1$};
        \draw[->, line width=1.8pt, AWSBlue] (0,0) -- (0,0.7) node[left=2pt, font=\small\bfseries] {$\sigma_2 \mathbf{e}_2$};
      \end{scope}

      % Arrow 2: Rotate into Output Space (U)
      \draw[->, ultra thick, UpGradSlate] (7.3, 0) -- (8.5, 0) node[midway, above=2pt, font=\small\bfseries\color{UpGradNight}] {$\mathbf{U}$};

      % Step 3: Rotated Ellipse in R^m (u1, u2 basis)
      \begin{scope}[shift={(10.6,0)}]
        \node[font=\small\bfseries\color{UpGradNight}] at (0, 2.3) {3. Output Space ($\mathbf{U}$)};
        \draw[thick, UpGradRule!70] (-2.0,0) -- (2.0,0);
        \draw[thick, UpGradRule!70] (0,-2.0) -- (0,2.0);
        \draw[thick, AWSEmerald, fill=AWSEmerald!12, rotate=30] (0,0) ellipse [x radius=1.9cm, y radius=0.7cm];
        \draw[->, line width=1.8pt, UpGradRed] (0,0) -- ({1.9*cos(30)}, {1.9*sin(30)}) node[above right=-1pt, font=\small\bfseries] {$\sigma_1 \mathbf{u}_1$};
        \draw[->, line width=1.8pt, AWSBlue] (0,0) -- ({-0.7*sin(30)}, {0.7*cos(30)}) node[above left=-1pt, font=\small\bfseries] {$\sigma_2 \mathbf{u}_2$};
      \end{scope}
    \end{tikzpicture}
  \end{center}
  \vspace{-0.10cm}
  \TakeawayBanner{Every real matrix $\mathbf{X} = \mathbf{U} \Sigma \mathbf{V}^T$: Orthogonal Rotation $\to$ Axis Stretch $\to$ Orthogonal Rotation.}
\end{frame}

% -----------------------------------------------------------------------------
% Frame 5: Production Checkpoint: PCA from Scratch via SVD
% -----------------------------------------------------------------------------
\begin{frame}[fragile]{Production Checkpoint: Principal Component Analysis via SVD}
  \begin{columns}[T,onlytextwidth]
    \begin{column}{0.48\textwidth}
      {\small\bfseries\color{UpGradNight}PCA Mathematical Mechanics}\par\vspace{0.08cm}
      For zero-centered data $\mathbf{X} \in \mathbb{R}^{N \times d}$:
      \begin{enumerate}\setlength{\itemsep}{0.20em}\small
        \item Compute SVD: $\mathbf{X} = \mathbf{U} \Sigma \mathbf{V}^T$.
        \item Right singular vectors $\mathbf{V} \in \mathbb{R}^{d \times d}$ are the principal loading directions.
        \item Project to $k$ dimensions: $\mathbf{Z} = \mathbf{X} \mathbf{V}_{:,:k} = \mathbf{U}_{:,:k} \Sigma_{:k,:k}$.
        \item Explained variance ratio: $\frac{\sigma_i^2}{\sum \sigma_j^2}$.
      \end{enumerate}
    \end{column}
    \hfill
    \begin{column}{0.48\textwidth}
      {\small\bfseries\color{UpGradRed}NumPy Implementation from Scratch}\par\vspace{0.06cm}
\begin{CodeBlock}{Python}
import numpy as np

def pca_from_scratch(X: np.ndarray, k: int):
    # 1. Center the features (zero mean)
    X_centered = X - np.mean(X, axis=0)
    
    # 2. Economy SVD decomposition
    U, S, Vt = np.linalg.svd(X_centered, full_matrices=False)
    
    # 3. Project data onto top-k components
    Z = X_centered @ Vt[:k].T  # shape: (N, k)
    
    # 4. Compute explained variance ratio
    explained_var = (S**2) / np.sum(S**2)
    return Z, Vt[:k], explained_var[:k]
\end{CodeBlock}
    \end{column}
  \end{columns}
\end{frame}

\end{document}
